\textbf{1. Kiến thức cần nhớ}
\begin{itemize}
\item Véctơ chỉ phương (VTCP) $\vec{u} \neq \vec{0}$: có giá song song hoặc trùng với đường thẳng.
\item Véctơ pháp tuyến (VTPT) $\vec{n} \neq \vec{0}$: có giá vuông góc với đường thẳng.
\item Mối liên hệ: Nếu đường thẳng có VTPT $\vec{n} = (a; b)$ thì có VTCP $\vec{u} = (-b; a)$ hoặc $(b; -a)$ và $\vec{n} \cdot \vec{u} = 0$.
\item Phương trình tổng quát (PTTQ) của đường thẳng đi qua $M(x_0; y_0)$ và có VTPT $\vec{n} = (a; b)$:
  $\boxed{a(x-x_0) + b(y-y_0) = 0} \Leftrightarrow ax + by + c = 0$.
\item Phương trình tham số (PTTS) của đường thẳng đi qua $M(x_0; y_0)$ và có VTCP $\vec{u} = (a; b)$:
  $\boxed{\begin{cases} x = x_0 + at \\ y = y_0 + bt \end{cases}} \quad (t \in \mathbb{R})$.
\item Khoảng cách từ điểm $M(x_0; y_0)$ đến đường thẳng $\Delta: ax+by+c=0$[cite: 3]:
  $\boxed{d(M, \Delta) = \frac{|ax_0+by_0+c|}{\sqrt{a^2+b^2}}}$.
\item Góc giữa hai đường thẳng $\Delta_1, \Delta_2$ có VTPT $\vec{n_1}, \vec{n_2}$[cite: 3]:
  $\cos \varphi = \frac{|\vec{n_1} \cdot \vec{n_2}|}{|\vec{n_1}| \cdot |\vec{n_2}|}$.
\end{itemize}

\textbf{2. Phương pháp giải các dạng}

%% DANG 0H9?3-1
\textbf{Dạng 1. Điểm, véctơ, hệ số góc của đường thẳng}
\begin{enumerate}
\item Từ PTTQ $ax+by+c=0$, suy ra VTPT $\vec{n}=(a;b)$ và VTCP $\vec{u}=(-b;a)$.
\item Từ PTTS, suy ra VTCP $\vec{u}=(a;b)$, điểm đi qua $M(x_0; y_0)$ và hệ số góc $k = \frac{b}{a}$ (nếu $a \neq 0$).
\end{enumerate}
Lưu ý: Một đường thẳng có vô số VTPT và VTCP. Các VTPT (hoặc VTCP) đều cùng phương với nhau. Điểm thuộc đường thẳng khi thay tọa độ vào phương trình phải thỏa mãn (hai vế bằng nhau).

%% DANG 0H9?3-2
\textbf{Dạng 2. Phương trình đường thẳng}
\begin{enumerate}
\item Tìm tọa độ $1$ điểm $M(x_0; y_0)$ mà đường thẳng đi qua.
\item Tìm tọa độ VTPT $\vec{n}$ (để viết PTTQ) hoặc VTCP $\vec{u}$ (để viết PTTS) dựa vào tính chất song song, vuông góc.
\item Ráp vào công thức PTTQ hoặc PTTS tương ứng.
\end{enumerate}
Lưu ý: $\Delta_1 \perp \Delta_2 \Rightarrow$ VTCP của đường này là VTPT của đường kia; $\Delta_1 \parallel \Delta_2 \Rightarrow$ hai đường thẳng dùng chung VTPT và VTCP.

%% DANG 0H9?3-3
\textbf{Dạng 3. Vị trí tương đối giữa hai đường thẳng}
\begin{enumerate}
\item Xét hệ phương trình gồm $2$ phương trình của $2$ đường thẳng.
\item Giải hệ bằng máy tính: Nếu hệ có $1$ nghiệm $\Rightarrow$ Cắt nhau; vô nghiệm $\Rightarrow$ Song song; vô số nghiệm $\Rightarrow$ Trùng nhau.
\item Hoặc xét tỉ số từ PTTQ[cite: 3]: $\frac{a_1}{a_2} \neq \frac{b_1}{b_2}$ (cắt); $\frac{a_1}{a_2} = \frac{b_1}{b_2} \neq \frac{c_1}{c_2}$ (song song); bằng nhau hết (trùng).
\end{enumerate}
Lưu ý: Điều kiện xét tỉ số chỉ dùng khi các hệ số ở mẫu khác $0$. Nếu có hệ số bằng $0$, bắt buộc phải giải hệ phương trình giao điểm.

%% DANG 0H9?3-4
\textbf{Dạng 4. Bài toán về góc giữa hai đường thẳng}
\begin{enumerate}
\item Xác định đúng cặp VTPT $\vec{n_1}, \vec{n_2}$ hoặc cặp VTCP $\vec{u_1}, \vec{u_2}$ của hai đường thẳng.
\item Áp dụng công thức tính $\cos \varphi = \frac{|\vec{n_1} \cdot \vec{n_2}|}{|\vec{n_1}| \cdot |\vec{n_2}|}$ (tử số có trị tuyệt đối).
\item Bấm máy tính (Shift Cos) để tìm số đo góc $\varphi$.
\end{enumerate}
Lưu ý: Góc giữa hai đường thẳng luôn $\le 90^\circ$[cite: 3], khác với góc giữa hai véctơ có thể là góc tù. Nếu tính ra cosin âm tức là quên trị tuyệt đối ở tử số.

%% DANG 0H9?3-5
\textbf{Dạng 5. Bài toán về khoảng cách}
\begin{enumerate}
\item Phải đưa phương trình đường thẳng về đúng dạng tổng quát: $ax+by+c=0$.
\item Xác định tọa độ điểm $M(x_0; y_0)$ cần tính khoảng cách.
\item Thay vào công thức $d = \frac{|ax_0+by_0+c|}{\sqrt{a^2+b^2}}$ và tính kết quả[cite: 3].
\end{enumerate}
Lưu ý: Để tính khoảng cách giữa hai đường thẳng song song, ta lấy một điểm bất kỳ trên đường thẳng này và tính khoảng cách từ điểm đó đến đường thẳng kia.

%% DANG 0H9?3-6
\textbf{Dạng 6. Bài toán tìm điểm}
\begin{enumerate}
\item Tham số hóa tọa độ điểm cần tìm. Nếu điểm thuộc đường thẳng có PTTS thì điểm có tọa độ dạng $(x_0+at; y_0+bt)$.
\item Dùng dữ kiện bài toán (cách đều, khoảng cách, diện tích...) lập phương trình ẩn $t$.
\item Giải phương trình tìm $t$, thay ngược lại để kết luận tọa độ điểm.
\end{enumerate}
Lưu ý: Nếu điểm thuộc đường thẳng dạng PTTQ $ax+by+c=0$, có thể rút $y$ theo $x$ để gọi tọa độ, ví dụ điểm có hoành độ $m$ thì tung độ là $\frac{-am-c}{b}$.

%% DANG 0H9?3-7
\textbf{Dạng 7. Bài toán dùng cho tam giác, tứ giác}
\begin{enumerate}
\item Dịch các tính chất hình học sang tọa độ: Đường cao thì vuông góc với cạnh đối diện, trung tuyến thì đi qua trung điểm...
\item Tìm tọa độ đỉnh bằng cách giải hệ phương trình giao điểm của hai đường thẳng cắt nhau tại đỉnh đó.
\item Kết hợp các công thức chu vi, diện tích, khoảng cách để xử lý yêu cầu bài toán.
\end{enumerate}
Lưu ý: Nhớ tính chất các điểm đặc biệt: Trực tâm là giao $2$ đường cao, trọng tâm là giao $2$ trung tuyến, tâm ngoại tiếp là giao $2$ trung trực.

%% DANG 0H9?3-8
\textbf{Dạng 8. Bài toán thực tế, PP tọa độ hóa}
\begin{enumerate}
\item Gắn hệ trục tọa độ $Oxy$ vào mô hình thực tế, ưu tiên chọn gốc tọa độ $O$ tại các vị trí đặc biệt (tâm, góc vuông).
\item Chuyển đổi các số liệu độ dài, khoảng cách thực tế thành tọa độ điểm trên mặt phẳng.
\item Sử dụng các công thức phương trình đường thẳng, khoảng cách, góc để giải quyết bài toán.
\end{enumerate}
Lưu ý: Phải thống nhất một đơn vị đo chiều dài thực tế (ví dụ km) tương ứng với $1$ đơn vị độ dài trên hệ trục tọa độ $Oxy$.

%% DANG 0H9?3-9
\textbf{Dạng 9. Bài toán có dùng PT chính tắc}
\begin{enumerate}
\item Xác định điểm đi qua $M(x_0; y_0)$ và VTCP $\vec{u}=(a; b)$.
\item Kiểm tra điều kiện $a \neq 0$ và $b \neq 0$.
\item Viết phương trình chính tắc: $\frac{x-x_0}{a} = \frac{y-y_0}{b}$.
\end{enumerate}
Lưu ý: Nếu VTCP có $1$ tọa độ bằng $0$ (tức là $a=0$ hoặc $b=0$) thì đường thẳng KHÔNG có phương trình chính tắc, khi đó bắt buộc phải dùng PTTQ hoặc PTTS.
